\UnnumberedChapter{LỜI MỞ ĐẦU}

Hạnh phúc là một trong những chỉ số quan trọng phản ánh chất lượng cuộc sống của mỗi quốc gia. Bên cạnh các tiêu chí kinh tế truyền thống, mức độ hạnh phúc còn chịu ảnh hưởng bởi hỗ trợ xã hội, tuổi thọ khỏe mạnh, quyền tự do lựa chọn, sự hào phóng và nhận thức về tham nhũng. World Happiness Report cung cấp một bộ dữ liệu có cấu trúc để khảo sát và so sánh các yếu tố này giữa các quốc gia qua nhiều năm.

Trong học phần Trực quan hóa dữ liệu, Nhóm 11 lựa chọn thực hiện đề tài \textbf{“Hệ thống trực quan hóa dữ liệu World Happiness Report”}. Nhóm xây dựng dashboard bằng Python, Streamlit, Pandas và Plotly nhằm chuyển dữ liệu thô thành các biểu đồ, bảng thống kê và nhận xét dễ quan sát. Hệ thống cho phép người dùng xem dữ liệu theo từng năm, theo dõi xu hướng giữa các năm, phân tích theo khu vực, đánh giá tác động của các yếu tố hạnh phúc và so sánh trực tiếp các quốc gia.

Mục tiêu của đề tài là tổ chức dữ liệu hợp lý, tạo các biểu diễn trực quan phù hợp và rút ra những nhận xét có cơ sở về mức độ hạnh phúc của các quốc gia cũng như mối liên hệ giữa điểm hạnh phúc với các chỉ số kinh tế -- xã hội. Phạm vi dữ liệu của đề tài gồm báo cáo World Happiness Report giai đoạn 2015--2019.

Bài tập lớn được trình bày qua ba chương chính. Chương 1 giới thiệu đề tài, dữ liệu, mục tiêu và các công cụ sử dụng. Chương 2 mô tả quy trình tiền xử lý dữ liệu và xây dựng hệ thống trực quan. Chương 3 trình bày kết quả phân tích qua các trang dashboard và các nhận xét rút ra từ dữ liệu.
